% !TEX root = userguide.tex

\chapter{Add-ons}
\pgst{tfemch}

\label{ch:addons}

\def\chaptertitle{Add-ons}
\def\headdate{11th September 2025}

In this chapter a short overview of the available add-ons is given in alphabetical order.

\section{bcf}
\label{sec:bcfaddon}

The add-on `bcf' contains routines for simulation of
stochastic fluid models in isolation or as part of a flow simulation using
Brownian configuration fields.

\subsection{Stochastic models}

Currently one module is available, and can be used as follows:
\begin{quote}
   \texttt{use stochastic\_models\_2D\_m}
\end{quote}
for solving stochastic fluid models in isolation using Brownian dynamics. The
models available are described in Appendix~\ref{chap:stochasticmodels}

The properties of the model are stored in a structure of type
\texttt{stmodel\_t}. The structure must first be constructed by calling the
routine \texttt{create\_stochastic\_model}. Then the material parameters must
be filled in the components \texttt{modulus}, \texttt{lambda}, \texttt{nonlin}
of the structure. The structure can be deleted using the generic routine
\texttt{delete}.

The Langevin equation is solved using Brownian dynamics. The numerical
parameters for this method should be stored in a structure of type
\texttt{stnumpar\_t}, with components \texttt{nfield} (number of
fields/realizations), \texttt{timeint} (type of time integration)
 and \texttt{tstep} (the time step).

With
\begin{quote}
   \texttt{getexample bcf}
\end{quote}
several examples, such as startup of shear and elongational flow are available.

\subsection{BCF flows}

The add-on `bcf' also contains routines for flows of
stochastic fluid models using Brownian configuration fields.
Currently only elements for 2D flows are available
 and can be used as follows:
\begin{quote}
   \texttt{use bcf\_elements\_m}
\end{quote}
The elements use DG for the convection term.

The parameters are
supplied to the elements using the structure \texttt{coefficients} of type
\texttt{coefficients\_t} (see Sec.~\ref{sec:coefficients}). The documentation
of the coefficients is not part of this manual, but is in the text files
\texttt{README\_coefficients\%i} and
\texttt{README\_coefficients\%r} in the directory
\texttt{<tfemroot>/src/addons/bcf}.

The theory and examples are described in Sec.~\ref{sec:bcfdg}.

\section{compressible\_fluid}
\label{sec:addon_compressible}

The add-on `compressible\_fluid' contains routines for adding (weak)
compressibility to fluids. It not suitable for gases.

The routines can be used as follows:
\begin{quote}
   \texttt{use compressible\_fluid\_m}
\end{quote}

Examples are described in Sec.~\ref{sec:compressible}.

\section{convection\_diffusion}
\label{sec:convdiff_elements}

The add-on `convection diffusion' contains routines for the Poisson equation
and the convection-diffusion equation. The routines for the Poisson equation
can be used as follows:
\begin{quote}
   \texttt{use poisson\_elements\_m}
\end{quote}
The routines for the convection-diffusion equation can be used as follows:
\begin{quote}
   \texttt{use convection\_diffusion\_elements\_m}
\end{quote}
Note, that the above statement also `uses' the poisson elements.
Both 2D and 3D elements can be used.

The parameters are
supplied to the elements using the structure \texttt{coefficients} of type
\texttt{coefficients\_t} (see Sec.~\ref{sec:coefficients}). The documentation
of the coefficients is not part of this manual, but is in the text files
\texttt{README\_coefficients\%i} and
\texttt{README\_coefficients\%r} in the directory
\texttt{<tfemroot>/src/addons/convection\_diffusion}.

The theory and examples are described in Sec.~\ref{sec:poissonweak},
Sec.~\ref{sec:diffusionweak} and subsequent sections.

\section{convection\_diffusion\_supg}
\label{sec:convdiff_elements_supg}

(This add-on has been developed by Giovanna Grosso, Caroline Balemans and Nick Jaensson of the TU/e)

The add-on `convection\_diffusion\_supg' contains routines for solving
the convection-diffusion equation. SUPG is optionally applied to stabilize the elements
for large P\'eclet numbers. Note, that the add-on is separate from the add-on
described in Sec.~\ref{sec:convdiff_elements}, having incompatible definitions for parameters in
the structure \texttt{coefficients}. Also, this add-on depends on the Stokes (Sec.~\ref{sec:stokesaddon}) and
supg utils (Sec.~\ref{sec:supg_utils}) add-ons.

The routines for the convection-diffusion equation can be used as follows:
\begin{quote}
   \texttt{use convection\_diffusion\_supg\_elements\_m}
\end{quote}
Both 2D and 3D elements can be used.

The parameters are
supplied to the elements using the structure \texttt{coefficients} of type
\texttt{coefficients\_t} (see Sec.~\ref{sec:coefficients}). The documentation
of the coefficients is not part of this manual, but is in the text files
\texttt{README\_coefficients\%i} and
\texttt{README\_coefficients\%r} in the directory
\texttt{<tfemroot>/src/addons/convection\_diffusion\_supg}.

The theory and examples are described in
Sec.~\ref{sec:convdiffsupg} and subsequent sections.

\section{dfm}
\label{sec:deformationfields}

The add-on deformation fields method (\texttt{dfm}) contains routines for solving
deformation fields.

The routines for the deformation fields method can be used as follows:
\begin{quote}
   \texttt{use dfm\_elements\_m}
\end{quote}
Both 2D and 3D elements can be used.

The parameters are
supplied to the elements using the structure \texttt{coefficients} of type
\texttt{coefficients\_t} (see Sec.~\ref{sec:coefficients}). The documentation
of the coefficients is not part of this manual, but is in the text files
\texttt{README\_coefficients\%i} and
\texttt{README\_coefficients\%r} in the directory
\texttt{<tfemroot>/src/addons/dfm}.

The theory and examples are described in
Sec.~\ref{sec:dfm} and subsequent sections.

\section{diffuse\_interface}
\label{sec:diffuse_interface}

The add-on `diffuse interface' contains routines for the diffuse interface
model.
The add-on can be used as follows:
\begin{quote}
   \texttt{use diffuse\_interface\_elements\_m}
\end{quote}
Both 2D and 3D elements can be used.
The module uses the module \texttt{stokes\_elements\_m}
 (see Sec.~\ref{sec:stokesaddon}).

The parameters are
supplied to the elements using the structure \texttt{coefficients} of type
\texttt{coefficients\_t} (see Sec.~\ref{sec:coefficients}). The documentation
of the coefficients is not part of this manual, but is in the text files
\texttt{README\_coefficients\%i} and
\texttt{README\_coefficients\%r} in the directory
\texttt{<tfemroot>/src/addons/diffuse\_interface}.

An example is described in Sec~\ref{sec:couettedi}.

\section{dummylibs}
\label{sec:dummylibs}

The add-on dummylibs can been used to provide dummy routines for missing
external libraries.

\section{elastic}
\label{sec:elastic}

The add-on `elastic' contains routines for nonlinear elastic solids. The routines
can be used as follows:
\begin{quote}
   \texttt{use elastic\_elements\_m}
\end{quote}
Both 2D and 3D elements can be used.

The parameters are
supplied to the elements using the structure \texttt{coefficients} of type
\texttt{coefficients\_t} (see Sec.~\ref{sec:coefficients}). The documentation
of the coefficients is not part of this manual, but is in the text files
\texttt{README\_coefficients\%i} and
\texttt{README\_coefficients\%r} in the directory
\texttt{<tfemroot>/src/addons/elastic}.

The theory is described in Sec.~\ref{sec:elasticsolid}.
An example is described in Sec.~\ref{sec:solidbar}. The add-on is also
being used in the fluid-solid interaction example in Sec.~\ref{sec:fluidsolid}.

\section{feast}
\label{sec:feast}

The add-on feast provides an interface to FEAST: a set of routines for solving
eigenvalue problems. It is free software. See the file \texttt{README} in the
directory \texttt{src/addons/feast} for some
directions on using FEAST. At the time of writing the add-on feast only interfaces to the routine in FEAST for solving real generalized eigenvalue problems for sparse matrices.

Examples using FEAST are
\texttt{couette23.f90}, \texttt{couette24.f90}, \texttt{couette25.f90} and \texttt{couette26.f90} which can be obtained by typing:
\begin{quote}
   \texttt{getexample couette}
\end{quote}
See Sec.~\ref{sec:linearstabilityCouette} for more info.

\section{figplot}
\label{sec:figplot}

The add-on figplot provides graphics routines writing a Fig output file. Fig files can be read and
edited using the xfig program and translated to various other formats.

The add-on figplot can be used by putting
\begin{quote}
   \texttt{use figplot\_m}
\end{quote}
in your program.

The add-on figplot provides several routines starting with \texttt{plot\_} that
are useful for producing 2D and projected 3D plots on any system
without the need for an external commercial plot program.
Note that figplot provides only basic plotting facilities
and is not intended to replace a full-blown external graphics program like
Tecplot or Paraview. In particular it only provides 2D plotting and projected 3D
plotting, i.e.\ the 3D plots are just simple projections on a 2D plane without
hidden line removal\footnote{By putting different graphical objects into
layers in the fig file a poor man's hidden line removal can be obtained.}.

Current possibilities are:
\begin{itemize}
   \item Plotting of points, curves and surfaces
        (routine \texttt{plot\_points\_curves}).
   \item Plotting of a mesh (routine \texttt{plot\_mesh}).
   \item Plotting of objects (routine \texttt{plot\_objects}).
   \item Plotting of data using color fill between contours (routine
         \texttt{plot\_color\_fill}).
   \item Plotting of data using colored contour lines (routine
         \texttt{plot\_color\_contour}).
   \item Plotting of two or three components with arrows (routine
         \texttt{plot\_vector}).
\end{itemize}
In 3D plotting is basically only possible on the specified surfaces.
Plotting options are set using a structure of type \texttt{plot\_options\_t}.
Directly setting plot options in the structure
\texttt{plot\_options} usually requires
a lot of lines of code. Also if more than one plot routine is called
it is more easy to restart from the default values.
For these reasons a helper routine (\texttt{set\_plot\_options}) has been
written which sets the structure to default value and the structure values can
be set by heading parameters. For example
\begin{listing}{1}

  call set_plot_options ( plot_options, numlevels=21 )

\end{listing}
Note, that \texttt{plot\_options} is set to default values first, before setting
the new values in the heading. Use \texttt{keep=.true.} in the heading to skip
the setting to default values.

A basic assumption underlies the \texttt{plot\_color\_xxx} and
\texttt{plot\_vector} routines:
\begin{quote}
   \emph{All nodes must contain data of the quantity to be plotted.}
\end{quote}
If a quantity violates this, a new derived vector should be created (using the
routine \texttt{derive\_vector}) that defines the quantity in all nodal points.
Examples are: pressures in velocity/pressure elements, a seven-node triangle
where the velocities are defined in the six boundary nodes and
the discontinuous stresses in the center node (node 7). In the latter example
both the velocities and stresses need derived vectors for plotting.

Each \texttt{plot\_xxx} routine produces a file written in the Fig-format as
used by the program \texttt{xfig}. The file should be given the extension
\texttt{.fig} and can be edited using \texttt{xfig}.
It can also be converted to other formats using the program \texttt{fig2dev}.

Two shell scripts have been written to ease the converting to the \texttt{eps}
format and to view the fig-files using ghostscript:
\begin{quote}
   \texttt{fig2eps [figfile(s)]}
\end{quote}
will generate eps-files of the figfiles given or all figfiles if no files are
given.
\begin{quote}
   \texttt{viewfig [figfile(s)]}
\end{quote}
shows on the screen the figfiles given or all figfiles if no files are
given. It uses ghostscript/ghostview.

Several shell scripts are available that convert fig-files to images.
For example
\begin{quote}
   \texttt{fig2tif [figfile(s)]}
\end{quote}
will generate tif-files of the figfiles given or all figfiles if no file are
given. Likewise, \texttt{fig2pcx}, \texttt{fig2gif} and \texttt{fig2png} are
available.

Many example files use figplot for plotting.

\section{generalized\_stokes}
\label{sec:generalized_stokes}

The add-on `generalized\_stokes' contains routines for
\begin{enumerate}
 \item The generalized Stokes equation, which is defined here as
a Stokes equation with a varying viscosity coefficient $\eta$. The varying $\eta$
coefficient needs to be supplied by the user and can
be specified by a function, a vector defined in the nodes or an element vector
defined per element (elvector).
 \item The flow of a generalized Newtonian fluid model, where the viscosity function $\eta(\dot\gamma)$ is
       according to
       \begin{alignat*}{3}
   \eta(\dot\gamma) &= m\dfrac{1}{\dot\gamma^{1-n}} & \quad&\text{Power-law}\\
   \eta(\dot\gamma) &= \eta_\infty + (\eta_0-\eta_\infty)
\dfrac{1}{\big(1+(\lambda\dot\gamma)^2\big)^{(1-n)/2}} &
          \quad&\text{Carreau}\\
   \eta(\dot\gamma) &= \eta_\infty + (\eta_0-\eta_\infty)
\dfrac{1}{\big(1+(\lambda\dot\gamma)^a\big)^{(1-n)/a}} &
          \quad&\text{Carreau-Yasuda}
\end{alignat*}
  with $\dot\gamma=\sqrt{2\ten D:\ten D}$.

The viscosity function $\eta(\dot\gamma)$ can be augmented with a yield function
 according to Papanastasiou \cite{Papanastasiou87}:
\[
     \eta(\dot\gamma) = \bar\eta(\dot\gamma) +
       \dfrac{\tau_\text{y} \left( 1 - \exp(-r \dot\gamma) \right)}{\dot\gamma}
\]
where $\bar\eta(\dot\gamma)$ is the viscosity function given by either a constant value (Bingham), a Power-law, Carreau or Carreau-Yasuda model (see above), $\tau_\text{y}$ is the yield stress and $r$ is the regularization parameter.

 \item Fluctuating hydrodynamics as introduced by Landau and Lifshitz.
 \item Building the stress work vector per element for viscous fluids.
\end{enumerate}
The add-on can be used as follows:
\begin{quote}
   \texttt{use generalized\_stokes\_elements\_m}
\end{quote}
Both 2D and 3D elements can be used.
The module uses the module \texttt{stokes\_elements\_m}
 (see Sec.~\ref{sec:stokesaddon}).

The parameters are
supplied to the elements using the structure \texttt{coefficients} of type
\texttt{coefficients\_t} (see Sec.~\ref{sec:coefficients}). The documentation
of the coefficients is not part of this manual, but is in the text files
\texttt{README\_coefficients\%i} and
\texttt{README\_coefficients\%r} in the directory
\texttt{<tfemroot>/src/addons/generalized\_stokes}.

Examples of varying viscosity are described in Sec~\ref{sec:genstokes}.
Examples of generalized Newtonian fluids are described in Sec~\ref{sec:gennewtonian}.
Examples of Brownian forces are described in Sec~\ref{sec:brownian}.

\section{grid\_deformation}
\label{sec:addon_grid_deformation}

(This add-on has been developed by Young Joon Choi of the TU/e.)

The add-on `grid deformation' contains routines for performing grid deformation
to change position and size of elements according to a monitor function.
The method is described in \cite{Wan07,Grajewski09} and used extensively in
the PhD thesis of Choi \cite{Choi11b}.

The routines can be used as follows:
\begin{quote}
   \texttt{use meshgen\_grid\_deformation\_m}\\
   \texttt{use deform\_mesh\_ma57\_m}
\end{quote}
where the first module is generic and the second module is specific for the
solver MA57 (see Sec.~\ref{sec:hslsparse}).

Examples are described in Sec.~\ref{sec:grid_deformation}.

\section{hsl}
\label{sec:hsl_addon}

The add-on hsl provides an interface to the direct solvers MA57 (symmetric) and MA41 from the HSL library \cite{HSL}. Furthermore there is an interface to the eigenvectors/value routine EA23 for diagonalizing $3\times3$ symmetric matrices.

See the file \texttt{README} in the directory \texttt{src/addons/hsl}.

\section{hsl\_extra}
\label{sec:hsl_extra_addon}

The add-on hsl\_extra adds additional interfaces to routines from the HSL library \cite{HSL}.
See the Sec.~\ref{sec:hsl_external} on how to obtain
the routines (requires a license, but academics can get a free personal license for each package separately).

The add-on includes the following interfaces \begin{itemize}
   \item Subroutine \texttt{condition\_number\_system\_mc75}: computes the
condition number of the system matrix using HSL routine MC75.
   \item Subroutine \texttt{condition\_number\_mc75}: computes the
condition number of a sparse matrix using HSL routine MC75.
   \item Subroutine \texttt{solve\_system\_mi20} solves a linear system using
an iterative solver combined with an algebraic multigrid (AMG) preconditioner
(HSL routine \texttt{hsl\_mi20}, see \cite{Doyle07}).
AMG preconditioners can achieve $\mathcal{O}(N)$ solvers for linear systems.
   \item Subroutine \texttt{solve\_system\_diag} solves a linear system using
an iterative solver combined with a diagonal (Jacobi) preconditioner.
    \item Special iterative solvers for the Stokes system, see Sec.~\ref{sec:iterativestokes}. The routines are
 \texttt{solve\_system1\_mi20}, \texttt{solve\_system2\_mi20} and \texttt{solve\_system3\_mi20}.
\end{itemize}
See the file \texttt{README} in the directory
\texttt{src/addons/hsl\_extra} for more information.

Examples for using  \texttt{solve\_system\_mi20}
 (iterative solve with amg
preconditioner) are
\texttt{poisson14.F90}, \texttt{poisson15.F90}, \texttt{poisson16.F90}
and \texttt{poisson17.F90}, which can be obtained by
typing:
\begin{quote}
   \texttt{getexample poisson}
\end{quote}
The problems are
the same as \texttt{poisson1.f90} (14) and \texttt{poisson5.f90} (15-17),
respectively,
but modified for using \texttt{solve\_system\_mi20}.
Use \texttt{diff} to see the differences.

Examples for using \texttt{solve\_system\_diag} are \texttt{poisson18.F90} and
 \texttt{poisson19.F90}.

\section{inertia}
\label{sec:inertia}

(This add-on has been developed by Michiel Baltussen of the TU/e).

The add-on `inertia' contains element routines for the inertia terms
in the momentum balance.
The element routines can be used as follows:
\begin{quote}
   \texttt{use inertia\_elements\_m}
\end{quote}
Both 2D and 3D elements can be used.
The module uses the module \texttt{stokes\_elements\_m}
(see Sec.~\ref{sec:stokesaddon}).

The parameters are
supplied to the elements using the structure \texttt{coefficients} of type
\texttt{coefficients\_t} (see Sec.~\ref{sec:coefficients}). The documentation
of the coefficients is not part of this manual, but is in the text files
\texttt{README\_coefficients\%i} and
\texttt{README\_coefficients\%r} in the directory
\texttt{<tfemroot>/src/addons/inertia}.

Examples are described in Sec~\ref{sec:navierstokes}.

\section{interface\_tracking}
\label{sec:interfacetrackingaddon}

(This add-on has been developed by Massimiliano Villone of the University of
Naples.)

The add-on `interface\_tracking' contains routines for tracking
sharp interfaces described by a mesh.
The basic routines can be used as follows:
\begin{quote}
   \texttt{use interface\_tracking\_elements\_m}\\
   \texttt{use poisson\_gd\_ma57\_m}
\end{quote}
where the first module contains the elements for tracking and the second module
a routine for the (optional) tangential motion of the interface. If the curvature is needed, for example for setting
the monitor function, an additional module needs to be loaded:
\begin{quote}
   \texttt{use curvature\_gd\_ma57\_m}
\end{quote}

The parameters are
supplied to the elements using the structure \texttt{coefficients} of type
\texttt{coefficients\_t} (see Sec.~\ref{sec:coefficients}). The documentation
of the coefficients is not part of this manual, but is in the text files
\texttt{README\_coefficients\%i} and
\texttt{README\_coefficients\%r} in the directory
\texttt{<tfemroot>/src/addons/interface\_tracking}.

Examples are described in Sec.~\ref{sec:interface_tracking}.

\section{interfacial\_rheology}
\label{sec:int_rheo}
The add-on `interfacial rheology' provides boundary element routines for
including interfacial rheology through a Neumann boundary condition in 2D,
3D, and axisymmetric interfacial or boundary flow problems. Viscous, elastic
and Kelvin-Voigt interface constitutive models have been implemented.
Furthermore, routines are available for computing the nodal interfacial
stresses on a 3D surface mesh.
The elements can be used as follows:
\begin{quote}
   \texttt{use interfacial\_rheology\_elements\_2D\_curve\_m}
\end{quote}
for 2D or axisymmetric problems and:
\begin{quote}
   \texttt{use interfacial\_rheology\_elements\_3D\_surface\_m}
\end{quote}
for 3D problems.

The parameters are supplied to the elements using the structure
\texttt{coefficients} of type \texttt{coefficients\_t} (see
Sec.~\ref{sec:coefficients}). The documentation of the coefficients is not
part of this manual, but is in the text files \texttt{README\_coefficients\%i}
and \texttt{README\_coefficients\%r} in the directory
\texttt{<tfemroot>/src/addons/interfacial\_rheology}.

\section{io\_utils}
\label{sec:io}

The add-on io\_utils supplies
routines for reading/writing mesh, problem, coefficients, data to
          and/or from a file. Also the interfacing (read/write) to external
          packages is in this module.

The add-on io\_utils can be used by putting
\begin{quote}
   \texttt{use io\_utils\_m}
\end{quote}
in your program.

The add-on io\_utils provides several special I/O routines. It consists of
several submodules that can also be used separately:
\begin{listing}{17}
  use tfem_utils_m  ! I/O for core tfem structures etc.
  use sep_utils_m   ! I/O utils for sepran
  use gbt_utils_m   ! I/O utils for Gambit (Neutral file format)
  use gmsh_utils_m  ! I/O utils for Gmsh (msh and pos file format)
  use ideas_utils_m ! I/O utils for NX-Ideas (Universal file format)
  use vtk_utils_m   ! I/O utils for VTK (legacy vtk file format)
  use tec_utils_m   ! I/O utils for Tecplot
  use printtofile_m ! Printing of data in vector/sysvectors to a file
  use printinfo_m   ! Printing of info from mesh, problem,... to standard output
\end{listing}

\subsection{Writing and reading a mesh to a file}
\label{sec:iomesh}

The routine \texttt{write\_mesh} can be used to write
basic mesh information (points, curves, surfaces, topology, coordinates)
to a binary file.
The routine \texttt{read\_mesh} can be used to read this information again,
possibly in a different main program.\\
Example writing the mesh:
\begin{listing}{1}

  <generate mesh>

  call write_mesh ( mesh, filename='mesh.out' )

\end{listing}
Example reading the mesh:
\begin{listing}{1}

  call read_mesh ( mesh, filename='mesh.out' )

  call fill_mesh_parts ( mesh )

\end{listing}
Note that only basic information is written and that after reading the mesh the
routine \texttt{fill\_mesh\_parts} must be called to generate all information of
a mesh.

\subsection{Writing and reading a structure of type
     \texttt{input\_probdef\_t} to a file}

In order to use vectors and sysvectors in a separate program for
post processing, the structure \texttt{problem} of type \texttt{problem\_t}
is needed. Since it is much easier to regenerate the structure
\texttt{problem} than the write/read it to/from a file, two routines
are available that write and read the information in the structure
\texttt{input\_probdef} of type \texttt{input\_probdef\_t}.

The routine \texttt{write\_input\_probdef} can be used to write
a structure of type \texttt{input\_probdef\_t} to a binary
file.
The routine \texttt{read\_input\_probdef} can be used to read this
information again, possibly in a different main program.\\
Example writing \texttt{input\_probdef}:
\begin{listing}{1}

! problem definition

  call create_input_probdef ( mesh, input_probdef, nvec=4, nphysq=3 )

  <fill input_probdef: element degrees of freedom, physical quantities,
                       essential degrees of freedom, constraints >

  call write_input_probdef ( mesh, input_probdef, filename='probdef.out' )

  <compute solution>

! write solution to binary file

  open ( unit=13, file='data.out', form='unformatted' )
  write ( unit=13 ) sol%u
  close ( unit=13 )

\end{listing}
Example reading \texttt{input\_probdef}:
\begin{listing}{1}

! problem definition

  call read_input_probdef ( mesh, input_probdef, filename='probdef.out' )

  call problem_definition ( input_probdef, mesh, problem )

! read solution from binary file

  open ( unit=13, file='data.out', form='unformatted' )
  read ( unit=13 ) sol%u
  close ( unit=13 )

  <post-processing sol>

\end{listing}

\subsection{Writing and reading coefficients to a file}
\label{sec:iocoef}

The routine \texttt{write\_coefficients} can be used to write
a structure of type \texttt{coefficients\_t} to a binary
file.
The routine \texttt{read\_coefficients} can be used to read this
information again, possibly in a different main program.\\
Example writing the coefficients:
\begin{listing}{1}

  <create coefficients and fill it>

  call write_coefficients ( coefficients, filename='coefficients.out' )

\end{listing}
Example reading the coefficients:
\begin{listing}{1}

  call read_coefficients ( coefficients, filename='coefficients.out' )

\end{listing}

\subsection{Reading a \sep-generated mesh}
\label{sec:readsepmesh}

The routine \texttt{read\_sepran\_mesh} reads a mesh from the a \sep-generated
mesh file \texttt{meshoutput}. See the source file \texttt{sep\_utils.f90} in the directory
\texttt{src/addons/io\_utils} for more information.
See the \tfem\ legacy repository for some
examples using \sep-generated meshes.


\subsection{Reading a Gmsh generated mesh (msh file format)}
\label{sec:readgmshmesh}

The routine \texttt{read\_mesh\_gmsh} reads a mesh from a file using
the msh format\footnote{Both a mesh in the ASCII and the binary format is
supported in \tfem.}
as generated by Gmsh\footnote{Gmsh
(\href{http://www.geuz.org/gmsh}{http://www.geuz.org/gmsh})
 \cite{Geuzaine09}, is an open source high quality mesh generator.
It uses a geometry file quite similar to \sep, but also has a graphical
interface.}. Some comments:
\begin{itemize}
\item The meshes generated are three dimensional. Use the optional argument
\texttt{ndim=} to create lower-dimensional meshes in \tfem.
\item Points, curves and/or surfaces are quite naturally generated.
\item  Physical entities can be used to
define the curves/surfaces
 using the optional argument \texttt{physgeom=.true.}. The
default is \texttt{physgeom=.false.}, which generates curves/surfaces using the
geometrical entities defined by Gmsh.
\item Using the optional argument \texttt{domaintype} the user can influence
which elements will become geometries (curves and surfaces) and which elements become
element groups.
\item Multiple element shapes can be used.
\item Multiple element groups of the same
element shape can be defined using physical entities.
\item Meshes with high-order polynomial interpolation can be read. By default, linear and quadratic elements are translated into lower-order elements of \tfem. Use the optional argument \texttt{highorder=.true.} to translate even the linear and quadratic elements to high-order elements in \tfem\footnote{Note, that the nodal topology will be different.}.
\end{itemize}
Note, that if physical entities are defined with gmsh,
only gmsh-elements that are part of physical entities will be in the msh file.
This means, for example, that if you want points, curves and (groups of)
elements in a 2D mesh, the required
points, curves and surfaces in gmsh need to be included in physical entities.
See the source file \texttt{gmsh\_utils.f90} in the directory
\texttt{src/addons/io\_utils} for more information.

Various examples (e.g.\ \texttt{poisson23.f90} and \texttt{ale9.f90}) use
a Gmsh generated mesh.

See also the text file \texttt{README\_gmsh} in the directory
\texttt{examples} for more information on the convention used
in \tfem\ regarding gmsh.

\subsection{Reading a Gambit-generated mesh}
\label{sec:readgbtmesh}

The routine \texttt{read\_mesh\_gambit} reads a mesh from a
file generated by Gambit\footnote{Gambit is the pre-processor
for the Fluent/Polyflow package.} in the neutral file format.
Note that:
\begin{itemize}
\item No points are generated.
\item
In order to generate nodesets, boundary conditions in nodes
need to be defined.
\item In order to generate curves and/or surfaces,
boundary conditions on element edges and/or faces need to be defined.
\item
Multiple element shapes can be used.
\item Multiple element groups of the same
element shape can be defined using material groups.
\end{itemize}
See the source file \texttt{gbt\_utils.f90} in the directory
\texttt{src/addons/io\_utils} for more information. No examples are available
in \tfem\ that use Gambit-generated meshes.

\subsection{Reading an NX-Ideas generated mesh (universal file format)}
\label{sec:readideasmesh}

(This module has been developed by Tae Gon Kang of the TU/e.)

The routine \texttt{read\_mesh\_ideas} reads a mesh from a file using
the universal format as
generated by NX-Ideas\footnote{NX-Ideas is part
of the UGS package for CAD/CAE/CAM.}.
See the source file \texttt{ideas\_utils.f90} in the directory
\texttt{src/addons/io\_utils} for more information.
See the \tfem\ legacy repository for some
examples using NX-Ideas-generated meshes.

\subsection{Writing to a VTK datafile}
\label{sec:vtk}

The Vizualization Toolkit\footnote{See
\href{http://www.vtk.org}{http://www.vtk.org}}
 (VTK)
is a free package for visualization of 2D
and 3D datasets. Only writing to the legacy format with the file
extension \texttt{.vtk} is available. Various free
scientific visualization programs can read
\texttt{.vtk} files, for example
Paraview\footnote{See
\href{http://www.paraview.org}{http://www.paraview.org}. An old userguide
manual can be downloaded
\href{http://www.paraview.org/files/v1.6/ParaViewUsersGuide.PDF}{here},
but a newer one must be bought.} or
VisIt\footnote{See \href{https://wci.llnl.gov/simulation/computer-codes/visit}{VisIt website}}.

The example file \texttt{stokes13.f90} writes data to the VTK format using the
two available routines \texttt{write\_scalar\_vtk} for writing a scalar
quantity, \texttt{write\_vector\_vtk} for writing a 2D or 3D vector quantity
and \texttt{write\_tensor\_vtk} for writing a 2D or 3D tensor
quantity\footnote{In the vtk output file a complete 3x3 tensor (9 components)
is written, whether it is 2D, 3D, symmetric or unsymmetric.
The file ends up being bigger than writing the individual components as a
scalar. The reason is the writing of zeros for the components that
are not actually there (2D) and the double writing of non-diagonal components
for a symmetric tensor.}.
The file can be obtained by typing at the
command line:
\begin{quote}
   \texttt{getexample stokes}
\end{quote}
It is an adaptation of the problem in \texttt{stokes11.f90} (3D driven cavity).
The writing is done at the lines
\begin{listing}{157}

! write to a vtk file for further post-processing

  call write_scalar_vtk ( mesh, problem, vector=pressure, dataname='pressure', &
    filename='cavity.vtk' )

  call write_vector_vtk ( mesh, problem, filename='cavity.vtk', &
    dataname='velocity_vector', sysvector=sol, append=.true. )

  call write_scalar_vtk ( mesh, problem, vector=dudz, filename='cavity.vtk', &
    dataname='dudz', append=.true. )

\end{listing}
In the first call the coordinates of the points and topology of the elements
are written in addition to the data. Using \texttt{append=.true.} in subsequent
calls, data can be added to the file for more scalar and vector quantities using
the existing points and topology. In Fig.~\ref{fig:cavity_dudz} the scalar
quantity $\plderiv{u}{z}$ is plotted using Paraview.
\begin{figure}[htb!]
   \begin{center}
   \includegraphics[scale=0.5]{cavity_dudz}
   \end{center}
   \caption{Isosurface plot of $\plderiv{u}{z}$ for \texttt{stokes13}.}
   \label{fig:cavity_dudz}
\end{figure}

The routines \texttt{write\_vector\_vtk}
and \texttt{write\_tensor\_vtk} assume that the dimension of vectors and tensors is the same as the space dimension (taken from the mesh). To override this default the optional arguments
\texttt{assume2D} and \texttt{assume3D} can be used.
To write 2D vectors and tensors on a 1D or 3D domain add the argument \texttt{assume2D=.true.}\ in the call of the routines.
To write 3D vectors and tensors on a 1D or 2D domain add the argument \texttt{assume3D=.true.}\ in the call of the routines. Note, that alternatively the optional argument \texttt{degfd} can be used. For example, the argument \texttt{assume3D} is equivalent to \texttt{degfd=[1,2,3]} for vectors.

It is also possible to write \tfem-vectors of type \texttt{vector\_t} having
\texttt{elementwise=.true.}. In this case, the data is written as
cell data, where in the vtk-file the data is a constant (scalar/vector/tensor)
value per cell/element. Since, the amount of data per element
for a vector with \texttt{elementwise=.true.} is usually more, the first data
available will be used. For a scalar this is usually the value in node 1, for a
vector the three components of node 1 and for a tensor the six/nine components
of node 1. However, other data layouts are possible, for example giving a scalar
value to the three nodes of a triangle is enough to be used as a vector in 3D
space. If \emph{both} nodal point data and cell data needs to be written,
don't forget
to use the \texttt{switch=.true.} parameter when switching from one data
format to the other. For example:
\begin{verbatim}

  call write_scalar_vtk ( mesh, problem, filename='sol.vtk', dataname='u', &
    sysvector=sol )
  call write_scalar_vtk ( mesh, problem, filename='sol.vtk', dataname='dudx', &
    degfd=1, append=.true., switch=.true., vector=grad )
  call write_scalar_vtk ( mesh, problem, filename='sol.vtk', dataname='dudy', &
    degfd=2, append=.true., vector=grad )

\end{verbatim}
where the vector \texttt{grad} has been defined as having at least
two components per element (for example two components in node 1). Note the use
of \texttt{switch=.true.} in the first call when going from nodal to cell data.

Note, that only constants per element can be written. In order to plot
higher-order functions per elements use \texttt{gmsh} (see
Sec.~\ref{sec:gmshwrite}).

Besides the routines \texttt{write\_scalar\_vtk},
\texttt{write\_vector\_vtk} and
\texttt{write\_tensor\_vtk} there are also the routines
\texttt{write\_mesh\_vtk} to write just the mesh with no data and \texttt{write\_geometry\_vtk}
to write the points/curves/surfaces of a mesh.

The default \emph{precision} for writing floating point data is determined by the value of \texttt{WRITE\_DOUBLE\_VTK} in the \texttt{limits\_m} module (see Sec.~\ref{sec:limits}). The value \texttt{WRITE\_DOUBLE\_VTK=.true.}\
  writes double precision (8 byte) and
 \texttt{WRITE\_DOUBLE\_VTK=.false.}\ writes single precision (4 byte).
 The value set in \texttt{limits\_m} is \texttt{WRITE\_DOUBLE\_VTK=.false.}\ (single precision), hence setting
\begin{verbatim}
   WRITE_DOUBLE_VTK = .true.
\end{verbatim}
before calls to the \texttt{write\_\ldots\_vtk} routines, will output all floating point data in double precision instead of single precision.
All \texttt{write\_\ldots\_vtk} routines have an optional \texttt{double} argument to impose the precision for individual calls. Note, that the output files will be 40--50\% larger for double precision as compared to single precision.

The default output \emph{format} for writing floating point data is determined by the value of \texttt{WRITE\_BINARY\_VTK} in the \texttt{limits\_m} module (see Sec.~\ref{sec:limits}). The value \texttt{WRITE\_BINARY\_VTK=.true.}\
  writes the data in a binary (internal) representation and
 \texttt{WRITE\_BINARY\_VTK=.false.}\ writes the data in a formatted (ASCII) representation.
The value set in \texttt{limits\_m} is \texttt{WRITE\_BINARY\_VTK=.true.}\ (binary output), hence setting
\begin{verbatim}
   WRITE_BINARY_VTK = .false.
\end{verbatim}
before calls to the \texttt{write\_\ldots\_vtk} routines, will output all data in the formatted (ASCII) format.
All \texttt{write\_\ldots\_vtk} routines have an optional \texttt{binary} argument to impose the format for individual calls. \red{WARNING: Make sure that \emph{all} data in a single file has the \emph{same format} (formatted or binary). Thus, when using the \texttt{append} argument, the value of the \texttt{binary} argument must be the same as for the first data written to the file.}
The advantages of a binary file are:
\begin{itemize}
 \item The data is identical to the internal representation when using double precision output (see above).
 \item The files are significantly smaller than formatted data files.
 \item Reading the file might be somewhat faster.
\end{itemize}
A disadvantage is, of course, that unformatted/binary files are unreadable by humans. Another possible, less common, disadvantage is the
incompatibility when transferring data files between computers that use a different internal binary format.

\subsection{Writing to a Tecplot datafile}
\label{sec:tec}

(This module has been developed by Tae Gon Kang of the TU/e.)

Tecplot is an easy to use commercial
scientific visualization package.
The file extension used is \texttt{.plt}. The following routines are available
\begin{verbatim}
   write_objects_tecplot
   write_scalar_tecplot
   write_mdata_tecplot
\end{verbatim}
to write data to a file. The usage is similar to the vtk-routines (see
Sec.~\ref{sec:vtk}) and the user is referred to the source code for the actual
heading of the routines.

The example file \texttt{stokes16.f90} writes data to the tecplot
format. The file can be obtained by typing at the
command line:
\begin{quote}
   \texttt{getexample stokes}
\end{quote}
It is an adaptation of the problem in \texttt{stokes9.f90}.

\subsection{Writing to a Gmsh datafile or mesh-only file}
\label{sec:gmshwrite}

Gmsh\footnote{Gmsh
(\href{http://www.geuz.org/gmsh}{http://www.geuz.org/gmsh})
 \cite{Geuzaine09}.} is primarily a high quality meshgenerator (see
Sec.~\ref{sec:readgmshmesh}), but can also be used to postprocess data files.
The datafiles are mesh files extended with data\footnote{Both the ASCII
and the binary format is supported in \tfem.}.
In \tfem\ we prefer to give these data files the file extension
\texttt{.mso} and reserve \texttt{.msh} for a file with a mesh only as written by gmsh.
A unique feature of gmsh (compared
to other external visualization tools) is the
ability to plot data that is discontinuous across element edges. This
can be achieved by writing a vector that has been defined with
\texttt{elementwise=.true.} (similar to
the internal graphics routines of figplot, see Sec.~\ref{sec:figplot}).

The example file \texttt{cylinder8.f90} (from the example
\texttt{confined\_cylinder}) writes data to the Gmsh format using the
routines \texttt{write\_scalar\_gmsh} for writing a scalar
quantity, \texttt{write\_vector\_gmsh} for writing a vector quantity
and \texttt{write\_tensor\_vtk} for writing a tensor
quantity\footnote{In the msh output file a complete 3x3 tensor (9 components)
is written, whether it is 2D, 3D, symmetric or unsymmetric.
The file ends up being bigger than writing the individual components as a
scalar. The reason is the writing of zeros for the components that
are not actually there (2D) and the double writing of non-diagonal components
for a symmetric tensor.}.

The routines \texttt{write\_vector\_gmsh}
and \texttt{write\_tensor\_gmsh} assume that the dimension of vectors and tensors is the same as the space dimension (taken from the mesh). To override this default the optional arguments
\texttt{assume2D} and \texttt{assume3D} can be used.
To write 2D vectors and tensors on a 1D or 3D domain add the argument \texttt{assume2D=.true.}\ in the call of the routines.
To write 3D vectors and tensors on a 1D or 2D domain add the argument \texttt{assume3D=.true.}\ in the call of the routines. Note, that alternatively the optional argument \texttt{degfd} can be used. For example, the argument \texttt{assume3D} is equivalent to \texttt{degfd=[1,2,3]} for vectors.

The file can be obtained by typing at the
command line:
\begin{quote}
   \texttt{getexample confined\_cylinder}
\end{quote}
It is an adaptation of the problem in \texttt{cylinder5.f90}. The procedure is
similar to the routines for writing VTK files (see Sec.~\ref{sec:vtk}).

Besides the routines \texttt{write\_scalar\_gmsh},
\texttt{write\_vector\_gmsh} and
\texttt{write\_tensor\_gmsh} there are also the routines
\texttt{write\_mesh\_gmsh}
 to write just the mesh  with no data and
\texttt{write\_geometries\_gmsh}
 to write just the geometries with no data. If the written \texttt{.msh} files
are used for
further computations, use \texttt{binary=.true.} in order
 to preserve the full accuracy
of the nodal coordinates.
Since points,
curves and surfaces are part of the written mesh, visualization with
gmsh can be used to identify these,
even for meshes made by other meshgenerators.
Writing only the geometries can be useful for remeshing with gmsh,
for example in an ALE computation.

It is also possible to write a scalar field to a so-called parsed format
(extension \texttt{.pos}) using the routine
 \texttt{write\_scalar\_gmsh\_parsed}.
Although the field can be viewed using gmsh, the
main application is to use the scalar data field as a ``background mesh'' for
setting target element sizes in adaptive mesh procedures. An example is given
in \texttt{adaptive\_mesh\_pos.f90}, which can be obtained by
\begin{quote}
   \texttt{getexample meshes}
\end{quote}

See the source file \texttt{gmsh\_utils.f90} in the directory
\texttt{src/addons/io\_utils} for more information.

\subsection{printtofile}
\label{sec:printtofile}

The module printtofile provides a routine \texttt{printtofile} to write the data
in vectors and sysvectors to a file for further processing, for example by
Matlab or \texttt{gnuplot}. Current possibilities are:
\begin{itemize}
   \item write data in the full domain or on curves/surfaces only.
         Setting global module variable
         \texttt{include\_node\_number=.false.} will also print the nodal
         number in the first column.
         Use one or more of the
         heading parameters \texttt{xmin}, \texttt{xmax},
        \texttt{ymin}, \texttt{ymax}, \texttt{zmin} and \texttt{zmax} to
         limit printing to a window.
   \item write 2D and 3D grid data. The grid data can only be written for
         regular grids.
         Note, that to avoid connecting some data points in gnuplot an
         empty line is written by default after one `line'.
         Set global module variable
         \texttt{blank\_line=.false.}, if this is not needed.
\end{itemize}
The output is a simple column structure with, optionally the nodal number,
the coordinates and the data. No particular format is used
(`\texttt{*}'-format).

\subsection{printinfo: printing some info of structures
 (mesh, problem, \ldots) to standard output}
\label{sec:printinfo}

The module printinfo provides a (generic) routine \texttt{printinfo} to write
info from structures to standard output. Currently only structures of type
\texttt{mesh\_t} and \texttt{problem\_t} are supported. For example
\begin{listing}{1}
   call printinfo ( mesh )
\end{listing}
will print some info of the mesh to standard output. A second optional
parameter is \texttt{printlevel} to print more or less than default:
\begin{listing}{1}
   call printinfo ( problem, printlevel=4 )
\end{listing}
will print the current maximum info on the problem.

\section{linear\_elastic}
\label{sec:linearelasticaddon}

The add-on `linear elastic' contains routines for solving linear elastic problems.
The elements can be used as follows:
\begin{quote}
   \texttt{use linear\_elastic\_elements\_m}
\end{quote}
Both 2D and 3D elements can be used. Elements using displacements only and mixed displacements-pressure elements are available.

The parameters are
supplied to the elements using the structure \texttt{coefficients} of type
\texttt{coefficients\_t} (see Sec.~\ref{sec:coefficients}). The documentation
of the coefficients is not part of this manual, but is in the text files
\texttt{README\_coefficients\%i} and
\texttt{README\_coefficients\%r} in the directory
\texttt{<tfemroot>/src/addons/linear\_elastic}.

Examples are described in Sec.~\ref{sec:linear_elastic}.

\section{material\_data}
\label{sec:materialdata}

The add-on \texttt{material\_data} provides the data for various materials. Currently, it gives the EGP model data for the materials PMMA, iPP, PS and PLLA in \cite{van_breemen_rate_2012}. The filling routines in the module become available after invoking:
\begin{quote}
   \texttt{use material\_data\_m}
\end{quote}
There are two variants: one that fills the data in a structure of type \texttt{vemodel\_t} and one that fills the data in a structure of type \texttt{coefficients\_t}. In the first case, the structure of type \texttt{vemodel\_t} needs to be created and filled, except for the filling of the actual material data. 
For example
\begin{listing}{1}
   call fill_data ( vemodel, 'PMMA' )
\end{listing}
will fill the data of the PMMA model in the structure \texttt{vemodel}. 
In the second case, the structure of type \texttt{coefficients\_t} needs to have been created and filled at least with the \emph{integer} values.
The example
\begin{listing}{1}
   call fill_data ( coefficients, 'iPP' )
\end{listing}
will fill the data of the iPP model in the structure \texttt{coefficients}. An optional argument can give the filled
data back in a more accessible format using a structure of type \texttt{vemodel\_t}:
\begin{listing}{1}
   call fill_data ( coefficients, 'iPP', vemodel )
\end{listing}
will fill the data of the iPP model in the structure \texttt{coefficients} and the structure \texttt{vemodel} of type \texttt{vemodel\_t}. However, the latter is only for easy access to the filled data and not for general use, since other data might not be correct.

\section{metis5}

\label{sec:metis5}
(The add-on for \metis-5 has been developed by Young Joon Choi).

The add-on \texttt{metis5} provides an interface to the library \metis-5, for renumbering to
reduce fill-in in sparse direct matrix solvers and also supports 64bit integers. The interface becomes available after invoking
\begin{quote}
   \texttt{use metis5\_m}
\end{quote}
\metis-5 is usually available as a precompiled library or requires a separate package download. See the file \texttt{README} in
the directory \texttt{src/addons/metis5}.

Most examples using sparskit (see Sec.~\ref{sec:sksolve}) or
mkl\_gmres (see Sec.~\ref{sec:mklgmres})
use \mbox{\metis-5}.

\section{mkl\_gmres}
\label{sec:mklgmres}

(This add-on has been developed by Gaetano d'Avino of the University of
Naples.)

The add-on mkl\_gmres provides an interface to the iterative solvers in MKL.
See the file \texttt{README} in the directory \texttt{src/addons/mkl\_gmres}
for some directions on using mkl\_gmres. Since the use is very similar to
the module \texttt{sk\_solve} it is recommended to read
Sec.~\ref{sec:sksolve} as well.

An example for using mkl\_gmres is in
\texttt{sphere37.F90}, which can be obtained by typing:
\begin{quote}
   \texttt{getexample sphere}
\end{quote}
The problem is
the same as in \texttt{sphere25.F90} (sparskit), but modified for
mkl\_gmres. Use \texttt{diff} to see the differences.

\section{mumps}
\label{sec:mumps}

The add-on mumps provides an interface to MUMPS: a freely available
parallel direct solver using MPI, but can also be used as an alternative
for MA41 by running a sequential version.
See the file \texttt{README} in the directory \texttt{src/addons/mumps}.

\section{pardiso}
\label{sec:pardiso}

The add-on pardiso provides an interface to PARDISO: a
parallel direct solver using OPENMP.
See the file \texttt{README} in the directory \texttt{src/addons/pardiso} for some
directions on using PARDISO.

An example for using PARDISO is in
\texttt{stokes43.F90} and \texttt{diffuse\_interface16.F90}.

\section{particle\_tracking}
\label{sec:particletrackingaddon}

The add-on `particle tracking' contains routines for tracking particles or
sharp interfaces consisting of particles. The number of particles can be
changed adaptively. The routines can be used as follows:
\begin{quote}
   \texttt{use particle\_tracking\_m}
\end{quote}
The implementation is based on the theory in \cite{Galaktionov00} (see also
\cite{Hulsen05a}).

This add-on uses the \texttt{odeint} ODE solver from Numerical Recipes (see Sec.~\ref{sec:nr}).

An example for using this add-on is in
\texttt{particle\_tracking1.F90},
which can be obtained by typing:
\begin{quote}
   \texttt{getexample particle\_tracking}
\end{quote}
An example \texttt{particle\_tracking2.F90} is available that is similar to \texttt{particle\_tracking1.F90}, but
now outputs data to files that can be read by Matlab to plot `filled blobs' as a function of time. A small Matlab script \texttt{plotparticle.m} is supplied as well.

\section{projection}
\label{sec:projection_addon}

The add-on `projection' contains routines for
\begin{enumerate}
   \item projection of a vector defined on one mesh to a vector defined on
         another mesh,
   \item projection of a vector to another vector, both defined on the same
         mesh, but having a different shape function.
\end{enumerate}

The routines can be used as follows:
\begin{quote}
   \texttt{use projection\_elements\_m}
\end{quote}

The parameters are
supplied to the elements using the structure \texttt{coefficients} of type
\texttt{coefficients\_t} (see Sec.~\ref{sec:coefficients}). The documentation
of the coefficients is not part of this manual, but is in the text files
\texttt{README\_coefficients\%i} and
\texttt{README\_coefficients\%r} in the directory
\texttt{<tfemroot>/src/addons/projection}.

Examples are described in Sec.~\ref{sec:projection}.

\section{shapefunc\_extra}
\label{sec:shapefunc_extra}

The add-on `shapefunc extra' contains routines for shape functions.
Currently two different modules are available, and
can be used as follows:
\begin{quote}
   \texttt{use shapefunc\_gauss\_m}
\end{quote}
for shape functions having the nodal points in the Gauss points and
\begin{quote}
   \texttt{use shapefunc\_hierarchic\_m}
\end{quote}
for hierarchic shape functions.

\section{sk\_solve}
\label{sec:sksolve}

The add-on sk\_solve provides interfaces to additional iterative solvers
from the Sparskit2 library. The external library Sparskit2 needs to be installed.
See the file \texttt{README} in the directory
\texttt{src/addons/sk\_solve} and Sec.~\ref{sec:sparskit2}.

Renumbering the nodes of the mesh for reduced fill-in
in sparse direct solvers has a positive effect on the memory used and/or the
convergence rate.
The add-on metis5 (see Sec.~\ref{sec:metis5}) can be used to renumber
the nodes, but a mesh made with an external meshgenerator may already have
been renumbered in a sufficient way.

Note, that for matrices with zeros on the diagonal (like in a
velocity-pressure system) a renumbering of the unknowns
such that the zero diagonal elements end up at the lower right block of the
system matrix, might help in getting a better preconditioner and faster convergence.
However, sometimes this has an opposite effect and
it is recommended to test it before using it for production runs.
An example for this and the use of the sparskit is in
\texttt{stokes12.F90}, which can be obtained by typing:
\begin{quote}
   \texttt{getexample stokes}
\end{quote}
The problem is
the same as in \texttt{stokes7.f90}, but modified for sparskit.
In the file we see
\begin{listing}{200}
! problem definition

  call create_input_probdef ( mesh, input_probdef, nvec=3, nphysq=2, &
    nphysqshifted=1 )

  input_probdef%vec_elementdof(1)%a =   &
      reshape ( [2,2,2,2,2,2,2,2,2,    &  ! velocity
                 1,0,1,0,1,0,1,0,0,    &  ! pressure
                 1,1,1,1,1,1,1,1,1 ], &  ! scalar, such as vorticity
                 [9,3] )

  input_probdef%physq = [1,2]
  input_probdef%physqshifted = [2]

\end{listing}
With the parameter \texttt{nphysqshifted} in the heading of
\texttt{create\_input\_probdef} we indicate the number of physical quantities
that need to be shifted towards the end of the vector (in this case only one:
the pressure). On line 212 the array \texttt{input\_probdef\%physqshifted}
is filled with the actual physical quantities involved in the shifting.

The module contains an input parameter \texttt{maxmvm} (in
the structure \texttt{solver\_options}) that must be set to the
maximum number of iterations (matrix-vector multiplies). If the number of
iterations is exceeded the program stops with an error message. The module
contains a logical \texttt{stoponmaxmvm} (also in
the structure \texttt{solver\_options}) to change this behavior: if
\texttt{stoponmaxmvm} is set to \texttt{.false.}, the solver still obeys the
maximum number of iterations, but the program continues. To detect whether in
that case the maximum number of iterations has been reached, there is a logical
output parameter \texttt{maxmvmreached} (in the heading of
the routine \texttt{solve\_system\_sk}) to test that. There parameters
are used in the following example of a
time-stepping scheme where the matrix changes every time step and we want to
keep the preconditioner as long as possible:
\begin{listing}{1}
  if ( step == 1 ) then
    solver_options%stoponmaxmvm = .true.
    call solve_system_sk ( sysmatrix, rhsd, sol, ilu, &
      solver_options=solver_options )
  else
    solver_options%stoponmaxmvm = .false.
    call solve_system_sk ( sysmatrix, rhsd, sol, ilu, initsol=.true., &
      solver_options=solver_options )
    if ( maxmvmreached ) then
      call delete_ilu_sk ( ilu )
      solver_options%stoponmaxmvm = .true.
      call solve_system_sk ( sysmatrix, rhsd, sol, ilu, initsol=.true., &
        solver_options=solver_options )
    end if
  end if
\end{listing}
In the first time step a preconditioner is built and kept. The solution is
required to be obtained within \texttt{maxmvm} iterations. In the following time
steps first an attempt is made to obtain the solution using the
old preconditioner. If that fails, the preconditioner is rebuild and a new
attempt is made with the new preconditioner. If it fails every time step, the
above procedure should not be used, of course. Note, that we use the solution
from the previous time step as an initial guess for the next.

In the example \texttt{sphere}:
\begin{quote}
   \texttt{getexample sphere}
\end{quote}
the file \texttt{sphere25.F90} is similar to \texttt{sphere36.f90} (see
Sec.~\ref{sec:sphere2}) but using an iterative solver from the Sparskit2
library.
Also
the file \texttt{sphere13.F90} is similar to \texttt{sphere3.f90} (see
Sec.~\ref{sec:sphereimplicit}) but using an iterative solver from the Sparskit2
library.

In Sparskit2, the 2-norm (Euclidian distance to the origin) of the residual is
used to check for convergence. Sparskit2 allows for the use of both a
relative criterion $\epsilon_\text{rel}$ and an absolute
criterion $\epsilon_\text{abs}$ and the iteration will be terminated if the
following condition is satisfied:
\[
 ||\col r_i|| < \epsilon_\text{rel}||\col r_0|| + \epsilon_\text{abs}
\]
where $||\col r_i||$ is the 2-norm of the residual at the $i$-th iteration
step. Using $\epsilon_\text{rel}$, the residual will always be normalized by
the initial residual of the iteration. However, $||\col r_0||$ depends on the
initial guess for the solution, which is not always
preferred. An alternative is to use the 2-norm of the right hand side, $||\col
f||$, as a normalizing factor (which is equal to the residual
if the null vector is tried as the solution). An example using this procedure
is in \texttt{ale14.f90}.

\section{stokes}
\label{sec:stokesaddon}

The add-on `stokes' contains routines for solving the Stokes equation.
The elements can be used as follows:
\begin{quote}
   \texttt{use stokes\_elements\_m}
\end{quote}
Both 2D and 3D elements can be used.

In this add-on also elements are available for building mixed Stokes problems
and computing the streamfunction in 2D.

The parameters are
supplied to the elements using the structure \texttt{coefficients} of type
\texttt{coefficients\_t} (see Sec.~\ref{sec:coefficients}). The documentation
of the coefficients is not part of this manual, but is in the text files
\texttt{README\_coefficients\%i} and
\texttt{README\_coefficients\%r} in the directory
\texttt{<tfemroot>/src/addons/stokes}.

Many examples make use of this add-on.

The elements also contain routines for computing drag forces. An example is
given in Sec.~\ref{sec:sphere} for Stokes flow around a sphere.

\section{structures}
\label{sec:structures_addon}

The add-on `structures' contains routines for solving mechanical truss and beam structures.

The routines can be used as follows:
\begin{quote}
   \texttt{use structures\_elements\_m}
\end{quote}

The parameters are
supplied to the elements using the structure \texttt{coefficients} of type
\texttt{coefficients\_t} (see Sec.~\ref{sec:coefficients}). The documentation
of the coefficients is not part of this manual, but is in the text files
\texttt{README\_coefficients\%i} and
\texttt{README\_coefficients\%r} in the directory
\texttt{<tfemroot>/src/addons/structures}.

Examples can be obtained by
\begin{quote}
   \texttt{getexample structures}
\end{quote}
but these are only briefly described in Sec.~\ref{sec:structures_examples}.

\section{supg\_utils}
\label{sec:supg_utils}

The add-on `SUPG utils' contains routines for use with the SUPG method.
This is an add-on used by some element routines, such as in the add-on
viscoelastic (see Sec.~\ref{sec:viscoelastic}).

\section{surface\_advection}
\label{sec:surfaceadvection}

The add-on `surface advection' contains routines for computing the surface
advection equation. This equation (the kinematic condition)
is used in free surface problems that can be described by a height function.
See \cite{Choi11c,Choi11b} and the references therein for the details.

The routines can be used as follows:
\begin{quote}
   \texttt{use surface\_advection\_elements\_m}
\end{quote}

The parameters are
supplied to the elements using the structure \texttt{coefficients} of type
\texttt{coefficients\_t} (see Sec.~\ref{sec:coefficients}). The documentation
of the coefficients is not part of this manual, but is in the text files
\texttt{README\_coefficients\%i} and
\texttt{README\_coefficients\%r} in the directory
\texttt{<tfemroot>/src/addons/surface\_advection}.

Examples are described in Secs.~\ref{sec:swellNewtonian} and
\ref{sec:swellviscoelastic}. Also some
test examples for a 2D height function (3D velocity vector) can be obtained by
\begin{quote}
   \texttt{getexample surface\_advection}
\end{quote}
but these are not further described in this manual.


\section{surface\_convection\_diffusion}
\label{sec:surfaceconvectiondiffusionaddon}

(This add-on has been developed by Christos Mitrias of the TU/e.)

The add-on `surface\_convection\_diffusion' contains routines for calculating the concentration
of a quantity on an interface.
The routines can be used as follows::
\begin{quote}
   \texttt{use surface\_convection\_diffusion\_m}\\
\end{quote}
The parameters are
supplied to the elements using the structure \texttt{coefficients} of type
\texttt{coefficients\_t} (see Sec.~\ref{sec:coefficients}). The documentation
of the coefficients is not part of this manual, but is in the text files
\texttt{README\_coefficients\%i} and
\texttt{README\_coefficients\%r} in the directory
\texttt{<tfemroot>/src/addons/surface\_convection\_diffusion}.

Examples are described in Sec.~\ref{sec:surface_convection_diffusion}.

\section{timer}
\label{sec:timer}

The add-on timer supplies the tic/toc timer routines.

The timer routines can be used after the following use statement:
\begin{quote}
   \texttt{use timer\_m}
\end{quote}

To time the CPU time of a section of the code use
\begin{listing}{1}

  call tic

  <code to time>

  call toc

\end{listing}
and the output will be something like
\begin{quote}
   \texttt{Elapsed CPU time = 1.7377 seconds}
\end{quote}
If \texttt{toc} is given a string argument, such as
\begin{listing}{1}

  call tic

  <code to time>

  call toc ( 'build_system' )

  <code to time>

  call toc ( 'solve_system' )

\end{listing}
the output of the last \texttt{toc} will be something like
\begin{quote}
\texttt{Elapsed CPU time = 2.12468 incremental and 3.02354 total seconds
           in solve\_system}
\end{quote}
where the incremental time is the time since the last call of \texttt{toc}
(or since the call of \texttt{tic} if it is the first \texttt{toc} called).
Note that the routine \texttt{tic} is like starting a stopwatch, whereas
\texttt{toc} is reading the stopwatch (not stopping it). To reset the stopwatch
to zero \texttt{tic} should be called again.

There is global logical variable \texttt{timer}
to turn off the printing of the CPU timings without removing the
\texttt{tic} and \texttt{toc} calls. Thus
\begin{listing}{1}

  timer = .false.

  call tic

  <code to time>

  call toc

\end{listing}
will not print timings.

CPU time is an indication of the amount of work of the CPU (cpu cycles). This does not include the time for I/O. Moreover, in parallel processing the CPU time (the work to be done) usually increases, whereas the real time (wall clock time) should go down. In that case, the routines \texttt{tic\_w} and \texttt{toc\_w}, which print the wall clock time, should be used to find the speedup. If you want to switch between \texttt{tic/toc} and \texttt{tic\_w/toc\_w}, change the use statement to
\begin{verbatim}
   use timer_m, ticl => tic, tocl => toc
\end{verbatim}
and the actual calls to calls to \texttt{ticl/tocl}. This will give the CPU time. Change the use statement to
\begin{verbatim}
   use timer_m, ticl => tic_w, tocl => toc_w
\end{verbatim}
to get the wall clock time.

\section{umfpack}
\label{sec:umfpack}

The add-on umfpack provides an interface to UMFPACK: a set of routines for solving sparse linear systems via LU factorization. It is free software and is part of the SuiteSparse library developed by Tim Davis of Texas A\&M University.
See the file \texttt{README} in the directory \texttt{src/addons/umfpack} for some
directions on using UMFPACK.

An example for using UMFPACK is in
\texttt{sphere33.F90}, which can be obtained by typing:
\begin{quote}
   \texttt{getexample sphere}
\end{quote}
The problem is
the same as in \texttt{sphere3.f90}, but modified for UMFPACK. Use
\texttt{diff} to see the differences.

\section{viscoelastic}
\label{sec:viscoelastic}

The add-on `viscoelastic' contains routines for flows of
viscoelastic and elastoviscoplastic fluids. The elements
can be used as follows:
\begin{quote}
   \texttt{use viscoelastic\_elements\_m}
\end{quote}
Both 2D and 3D elements can be used\footnote{The DG elements are
not available in 3D.}

The viscoelastic add-on depends on the module \texttt{stokes\_elements\_m}
 (see Sec.~\ref{sec:stokesaddon}).

Also available is the module
   \texttt{use devss\_elements\_m},
which can be used for building the DEVSS system matrices,

The add-on also uses (with \texttt{xx=2D} or \texttt{3D})
\begin{quote}
   \texttt{use viscoelastic\_models\_xx\_m}\\
   \texttt{use viscoelastic\_models\_xx\_log\_m}
\end{quote}
which are described in Sec.~\ref{sec:vemodels}.

The parameters are
supplied to the elements using the structure \texttt{coefficients} of type
\texttt{coefficients\_t} (see Sec.~\ref{sec:coefficients}). The documentation
of the coefficients is not part of this manual, but is in the text files
\texttt{README\_coefficients\%i} and
\texttt{README\_coefficients\%r} in the directory
\texttt{<tfemroot>/src/addons/viscoelastic}.

Examples are described in Secs~\ref{sec:driven_cavity}--\ref{sec:couette}.

The elements also contain routines for computing drag forces.

\section{viscoelastic\_fh}
\label{sec:viscoelastic_fh}

The add-on viscoelastic\_fh for `viscoelastic fluctuating hydrodynamics' contains elements to
add fluctuating hydrodynamics terms to the constitutive equation.

The element routines can be used as follows:
\begin{quote}
   \texttt{use viscoelastic\_fh\_m}
\end{quote}
Both 2D and 3D elements can be used.

This add-on uses the viscoelastic and generalized stokes add-ons.

Examples are not yet available.


\section{viscoelastic\_models}
\label{sec:vemodels}

The add-on `viscoelastic models' contains routines for simulation of
viscoelastic and elastoviscoplastic fluid models. It is almost self-contained and only uses some basic 
modules from \tfem, such as some math modules and modules for defining the precision. Therefore it can be used without much knowledge of \tfem\ to study material models.

Currently four different modules are available, and
can be used as follows:
\begin{quote}
   \texttt{use viscoelastic\_models\_2D\_m}
\end{quote}
for planar and axisymmetrical flow,
\begin{quote}
   \texttt{use viscoelastic\_models\_2D\_log\_m}
\end{quote}
for planar and axisymmetrical flow in the log formulation \cite{Hulsen05},
\begin{quote}
   \texttt{use viscoelastic\_models\_2D\_b\_m}
\end{quote}
for planar and axisymmetrical flow in the b-tensor formulation,
\begin{quote}
   \texttt{use viscoelastic\_models\_3D\_m}
\end{quote}
for 3D flow,
\begin{quote}
   \texttt{use viscoelastic\_models\_3D\_log\_m}
\end{quote}
for 3D flow in the log formulation \cite{Hulsen05}
\begin{quote}
   \texttt{use viscoelastic\_models\_3D\_b\_m}
\end{quote}
for 3D flow in the b-tensor formulation.

Note, not all models
available in the module \texttt{viscoelastic\_models\_2D\_m} are available in
the other modules. Also, the number of models implemented in the b-tensor formulation is limited. However, setting
\begin{quote}
   \texttt{USE\_CONFORMATION\_MODEL=.true.}
\end{quote}
in the module \texttt{limits\_m} (see Sec.~\ref{sec:limits}) will transform any model available in the c-tensor formulation to the b-tensor formulation. This means a new model needs to be implemented only in the c-tensor formulation in order to be used in the b-tensor formulation. A drawback is that it takes a little more CPU time. If the model computation is time-critical, it is recommended to implement the model in the b-tensor formulation directly.

The properties of the model are stored in a structure of type
\texttt{vemodel\_t}. The structure must first be constructed by calling the
routine \texttt{create\_viscoelastic\_model}. Then the material parameters must
be filled in the components \texttt{modulus}, \texttt{lambda}, \texttt{nonlin}
of the structure. For adapted $\lambda$ models, additional material parameters need to be filled in the component \texttt{alam}. The structure can be deleted using the generic routine
\texttt{delete}.

The models available are described in Appendix~\ref{chap:viscoelasticmodels}.

With
\begin{quote}
   \texttt{getexample viscoelastic\_models}
\end{quote}
several examples, such as startup of shear and elongational flow, steady state solution obtained using Newton-Raphson iteration and eigenvalues/eigenvectors of the linearized right-hand side of the models are available. See the file \texttt{README} for a description of each example.

\section{ziggurat}
\label{sec:ziggurat}

The add-on ziggurat provides a fast random number generator for a normal
distribution.


